% Yêu cầu 6 (bắt buộc nhóm B) — Ngân sách: ~2.5 trang
\chapter{Đối chiếu với TCVN 11930:2017}\label{chap:tcvn}

\section{Giới thiệu tiêu chuẩn}
TCVN 11930:2017~\cite{tcvn11930} quy định yêu cầu cơ bản về an toàn hệ thống thông tin theo năm cấp độ (Điều 5 đến Điều 9). Mỗi cấp độ gồm yêu cầu quản lý và yêu cầu kỹ thuật; yêu cầu kỹ thuật chia thành bốn nhóm: an toàn mạng, máy chủ, ứng dụng và dữ liệu. Nhóm đối chiếu với yêu cầu kỹ thuật cấp độ 3 (Điều 7) và ghi kèm điều khoản tương ứng của cấp độ 2 (Điều 6) khi có. Đáng chú ý, mục 7.2.3.6~c nêu đích danh các dạng tấn công mà đề tài trình diễn: \textit{``Có phương án bảo vệ ứng dụng chống lại những dạng tấn công phổ biến: SQL Injection, OS command injection, RFI, LFI, Xpath injection, XSS, CSRF''}.

\section{Bảng đối chiếu}
Cột \textit{Mức}: \textbf{Đ} đáp ứng, \textbf{MP} một phần, \textbf{K} không đáp ứng (cần thành phần khác).

\begin{footnotesize}
\renewcommand{\arraystretch}{1.1}
\begin{longtable}{@{}>{\raggedright\arraybackslash}p{3.7cm}>{\raggedright\arraybackslash}p{5.6cm}>{\raggedright\arraybackslash}p{3.6cm}c@{}}
\caption{Đối chiếu điều khoản TCVN 11930:2017 -- cấu hình -- bằng chứng}\label{tab:tcvn}\\
\toprule
\textbf{Điều khoản} & \textbf{Cấu hình đáp ứng} & \textbf{Bằng chứng} & \textbf{Mức} \\ \midrule
\endfirsthead
\toprule
\textbf{Điều khoản} & \textbf{Cấu hình đáp ứng} & \textbf{Bằng chứng} & \textbf{Mức} \\ \midrule
\endhead
\bottomrule
\endfoot
7.2.3.6~c -- Chống SQLi, OS command injection, LFI, RFI, XSS, \ldots
& \texttt{SecRuleEngine On}; CRS 4.30.0 nhóm 930 (LFI), 931 (RFI), 932 (RCE), 941 (XSS), 942 (SQLi); PL1, ngưỡng vào 5
& Bảng~\ref{tab:ketqua}, Hình~\ref{fig:sqli}
& MP \\
7.2.3.6~a (6.2.3.4) -- Kiểm tra dữ liệu đầu vào
& Pha 1--2 kiểm tra header, tham số, thân yêu cầu; nhóm 920 kiểm tra tuân thủ HTTP
& Listing~\ref{lst:audit-sqli}
& MP \\
7.2.3.6~b, d -- Kiểm tra dữ liệu đầu ra; kiểm soát thông báo lỗi
& \texttt{SecResponseBodyAccess On}; pha 4 với nhóm 950--959 (953100: lộ lỗi PHP), ngưỡng ra 4; trang 403 không lộ thông tin ứng dụng
& Mục 5.5 (luật 953100, 959100)
& MP \\
7.2.1.5~a (6.2.1.5~a) -- Phòng chống xâm nhập
& WAF dạng reverse proxy trước máy chủ web, chế độ Blocking
& Chương~\ref{chap:demo}
& Đ \\
7.2.1.5~b (6.2.1.5~b) -- Cập nhật dấu hiệu tấn công
& CRS cố định theo digest image; cập nhật bằng cách đổi digest rồi chạy lại bộ kiểm thử
& \texttt{docker-compose.yml}
& MP \\
7.2.1.5~c -- Năng lực đáp ứng quy mô
& Đo thông lượng ở PL1, PL4 (\texttt{scripts/bench.sh})
& Bảng~\ref{tab:bench}
& MP \\
7.2.1.4~a -- Nhật ký: thời gian, IP/cổng, hành động, cảnh báo
& Audit log JSON: \texttt{time\_stamp}, \texttt{client\_ip}, \texttt{client\_port}, \texttt{http\_code}, các luật khớp
& Listing~\ref{lst:audit-sqli}
& Đ \\
7.2.3.3~a (6.2.3.3~a) -- Nhật ký truy cập ứng dụng, lỗi
& Access log, error log của nginx; audit log
& \texttt{docker compose logs waf}
& MP \\
7.2.1.4~c, đ; 7.2.3.3~b, c -- Lưu tập trung, tối thiểu 03 tháng
& Chưa có: lab chỉ ghi log ra stdout của container
& --
& K \\
7.2.1.1~b -- Kiểm soát truy cập giữa các vùng; an toàn máy chủ CSDL
& Mạng \texttt{edge} và \texttt{backend} (\texttt{internal: true}); CSDL chỉ ở \texttt{backend}, không ra Internet
& \texttt{docker-compose.yml}
& MP \\
7.2.1.2~đ, 7.2.3.2~đ -- Giới hạn kết nối đồng thời
& CRS không có chức năng này
& --
& K \\
\end{longtable}
\end{footnotesize}

\section{Nhận xét mức độ đáp ứng}
WAF đáp ứng trực tiếp nhất mục 7.2.3.6~c: bốn trong bảy dạng tấn công được nêu (SQLi, OS command injection, LFI, XSS) đã bị chặn trong Chương~\ref{chap:demo}; CRS có nhóm 931 cho RFI nhưng nhóm chưa kiểm thử. Mức đánh giá chỉ là \textit{một phần} vì: CSRF không phát hiện được bằng chữ ký mà phải sửa ở ứng dụng (token chống CSRF); XPath injection chưa kiểm thử; và PL1 vẫn bỏ sót command injection dạng ngắn (mục~\ref{sec:hanche}). Các mục 7.2.3.6~a, b, d cũng chỉ đáp ứng một phần vì WAF là lớp bảo vệ chung, không thay được việc ứng dụng tự kiểm tra dữ liệu theo nghiệp vụ.

Audit log đã có đủ các trường mục 7.2.1.4~a yêu cầu, nhưng cấp độ 3 còn đòi lưu tập trung và giữ tối thiểu 03 tháng; lab mới ghi ra stdout nên cần đẩy log về một SIEM như Wazuh hoặc Loki/Grafana (yêu cầu 7 của đề tài). Giới hạn kết nối đồng thời nằm ngoài khả năng của CRS, có thể bổ sung bằng \texttt{limit\_conn}/\texttt{limit\_req} của nginx. Như vậy, một mình ModSecurity + CRS không đủ cho cấp độ 3.
